\section{Experimental Setup and Details}
\label{sec:appx:exp setup and details}
In this section, we provide comprehensive details for each application: layout-to-image generation, quantity-aware image generation, and aesthetic-preference image generation. We also include full experimental details.
\paragraph{Layout-to-Image Generation.}
This task involves placing user-specified objects within designated bounding boxes~\cite{li2023gligen, xiao2023rnb, phung2023attnrefocusing, lee2024groundit}. We evaluate performance on 50 randomly sampled cases from the HRS-Spatial~\cite{Eslam2023hrs} dataset. As the reward model, we use GroundingDINO~\cite{liu2023grounding}, and measure the alignment between predicted and target boxes using mean Intersection-over-Union (mIoU). For the held-out reward, we compute mIoU using a different object detector, Salience DETR~\cite{Hou2024saliencedetr}. 
\vspace{-0.3cm}
\paragraph{Quantity-Aware Image Generation.}
This task involves generating a user-specified object in a specified quantity~\cite{kang2025counting,binyamin2024count, zafar2024iterativecountgeneration}. We evaluate methods on a custom dataset constructed via GPT-4o~\cite{openai2024gpt4ocard}, comprising 20 object categories with randomly assigned counts up to 90, totaling 40 evaluation cases.
As the reward model, we used T2ICount~\cite{qian2025t2icount}, which takes a generated image and the corresponding text prompt as input and returns a density map. Summing over this density map yields a differentiable estimate of the object count $n_{\text{pred}}$. The reward is defined as the negative smooth L1 loss:
\begin{align}
    \nonumber r_{\text{count}}=
    \begin{cases}
        -0.5(n_{\text{pred}} - n_{\text{gt}})^2 & |n_{\text{pred}}-n_{\text{gt}}| < 1, \\
        -|n_{\text{pred}} - n_{\text{gt}}| + 0.5 & |n_{\text{pred}}-n_{\text{gt}}| \ge 1.
    \end{cases}
\end{align}
where $n_{\text{gt}}$ denotes the input object quantity. For the held-out reward, we used an alternative counting model, CountGD~\cite{AminiNaieni2024countgd}.
This model returns integer-valued object counts. We apply a confidence threshold of 0.3 and evaluate using mean absolute error (MAE) and counting accuracy, where a prediction is considered correct if $n_{\text{pred}}=n_{\text{gt}}$. 
\vspace{-0.3cm}
\paragraph{Aesthetic-Preference Image Generation}
This task involves generating visually appealing images. We evaluate performance using 45 prompts consisting of animal names, provided in~\cite{black2024ddpo}. As the reward model, we use the LAION Aesthetic Predictor V2~\cite{schuhmann2022aesthetic}, which estimates the aesthetic quality of an image and is commonly used in reward-alignment literature~\cite{kim2025DAS, tang2024finetuningdiffusionmodelsstochastic, uehara2024bridging, black2024ddpo}. 

\vspace{-0.3cm}
\paragraph{Common Held-Out Reward Models.}
For all applications, we additionally evaluate the generated images using widely adopted held-out reward models for image quality and text alignment. Specifically, we use ImageReward~\cite{xu2023imagereward}, fine-tuned on human feedback, and VQAScore~\cite{lin2024vqa}, which leverages a visual question answering (VQA) model. 
Both are based on vision-language models and assess how well the generated image aligns with the input text prompt.

\paragraph{Experimental Details.}
\vspace{-0.3cm}
We use FLUX-Schnell~\cite{BlackForestLabs:2024Flux} as the score-based generative model for our method and all baselines. Although FLUX is a flow-based model, SMC-based inference-time reward alignment can be applied by reformulating the generative ODE as an SDE~\cite{ma2024sitexploringflowdiffusionbased,kim2025inference}, as described in Sec.~\ref{sec:problem definition and background:score based gnerative models}. 
Further to ensure diversity between samples during SMC, we applied Variance Preserving (VP) interpolant conversion~\cite{kim2025inference, lipman2024flowmatchingguidecode}. 
Apart from FLUX, we additionally report quantitative and qualitative results using another score-based generative model, SANA-Sprint~\cite{chen2025sanasprint}, in Appendix~\ref{sec:appx:sana}. These results demonstrate that the our claims are not tied to a specific model architecture, and additionally highlighting the robustness of \methodname{}.\\

We use 25 denoising steps for all SMC-based methods and 50 for single-particle methods to compensate for their reduced exploration capacity. To ensure fair comparison, we fix the total number of function evaluations (NFE) across all SMC variants—1,000 for layout-to-image and quantity-aware generation tasks, and 500 for aesthetic-preference image generation. For aesthetic-preference image generation, we used half the NFE compared to other tasks because we found it sufficient for performance convergence.
For methods that sample from the posterior, half of the NFE is allocated to the initial sampling stage, resulting in 20 particles during SMC, while prior-based methods use all NFE for SMC with 40 particles. 
For the aesthetic-preference image generation, we use half the number of particles in both settings to reflect the halved NFE.
The ablation study comparing the performance on varying NFE allocation is provided in Appendix~\ref{sec:appx:ablation}.
For all experiments, we used NVIDIA A6000 GPUs with 48GB VRAM.
